% =====================================================================
% CAS PCBA Engineer — front matter
% Noetic Compass Press · 2026
% =====================================================================

\frontmatter
\pagestyle{plainbook}

% --- Half-title --------------------------------------------------------
\thispagestyle{empty}
\begin{center}
  \vspace*{2.5in}
  {\huge\bfseries CAS PCBA Engineer\par}
\end{center}
% \cleardoublepage, not \clearpage: the title page must open on a recto
% (right-hand) page, and the copyright must sit on the verso behind it.
\cleardoublepage

% --- Title page --------------------------------------------------------
\thispagestyle{empty}
\begin{center}
  \vspace*{1.1in}
  {\Huge\bfseries CAS PCBA Engineer\par}
  \vspace{0.55in}
  {\Large\itshape Designing Cognitive Circuits with the\\[0.10in] Cognitive Architecture System\par}
  \vspace{0.35in}
  {\large From First Resistor to Production Safety Board\par}
  \vspace{1.35in}
  {\Large Little Henderson\par}
  \vfill
  {\normalsize Noetic Compass Press\par}
  {\small 2026\par}
\end{center}
\clearpage

% --- Copyright ---------------------------------------------------------
\thispagestyle{empty}
\vspace*{\fill}
\begin{center}
\begin{minipage}{0.85\textwidth}
\small
\noindent
\textbf{CAS PCBA Engineer}\\
\textit{Designing Cognitive Circuits with the Cognitive Architecture System}\\[6pt]
Copyright \textcopyright{} 2026 Little Henderson. All rights reserved.\\[6pt]
No part of this publication may be reproduced, distributed, or transmitted in
any form or by any means, including photocopying, recording, or other electronic
or mechanical methods, without the prior written permission of the author, except
in the case of brief quotations embodied in critical reviews and certain other
noncommercial uses permitted by copyright law.\\[6pt]
First Edition, 2026. Published by Noetic Compass Press.\\[6pt]
Companion volume: \textit{Programming in CAS: A Practical Guide} (forthcoming).\\[6pt]
Framework: Cognitive Architecture System (CAS) v7.7.\\[6pt]
Base-framework intellectual property: USPTO Provisional Application
\#64/037,173. The Cognitive Junction Transistor (CJT), the Cognitive Circuit
Description Language (CCDL), and the Cognitive Circuit Board (CCB) reference
designs in this volume extend that foundation.\\[6pt]
ISBN: [pending publisher assignment]\\[6pt]
\textit{Product and company names mentioned in this book may be trademarks of
their respective owners, and are used here for identification and commentary
only.}\\[6pt]
\textit{Nothing in this book is a safety certification, a guarantee of model
behavior, or legal or medical advice. A Cognitive Circuit Board is an
engineering artifact: it does what its schematic, parameters, and test report
say it does --- no more, and no less.}
\end{minipage}
\end{center}
\vspace*{\fill}
\cleardoublepage

% --- Epigraph ----------------------------------------------------------
\thispagestyle{empty}
\vspace*{2.6in}
\begin{center}
\begin{minipage}{0.74\textwidth}
\itshape\centering
For everyone who ever read an AI safety incident report and thought:
we would never accept this from a bridge, a pacemaker, or a power
supply --- why do we accept it from a mind?
\end{minipage}
\end{center}
\clearpage

% --- Table of Contents -------------------------------------------------
{%
  \hypersetup{linkcolor=black}%
  \tableofcontents*
}
\clearpage

% --- List of Figures ---------------------------------------------------
{%
  \hypersetup{linkcolor=black}%
  \listoffigures
}
\clearpage

% --- Preface + How to Use This Book (unnumbered front matter) ----------
\input{front_prose.tex}
